% N4 --- Thevenin e Norton: 10 desafios
\blocodesafio

\begin{questao}[4]{}
Prove o teorema de Thévenin aplicando uma fonte de corrente de teste $I$ à porta e usando superposição.
\circuitoq{\begin{circuitikz}[american,scale=.82,transform shape]
\draw[dashed,cinza,rounded corners] (0,0) rectangle (3.4,3.8); \node[cinza] at(1.7,1.9){rede linear};
\draw (3.4,3.1)--(5,3.1) coordinate(a); \draw (3.4,.7)--(5,.7) coordinate(b);
\draw (b) to[isource,l_={$I$}] (a); \draw[<->,Vcol] (5.7,.9)--(5.7,2.9) node[midway,right] {$U$};
\end{circuitikz}}
\resposta{$U=V_{oc}-R_{Th}I$}
\resolucao{Por superposição, $U=U'+U''$. Com a fonte de teste zerada, a porta está aberta e $U'=V_{oc}$. Com as fontes internas zeradas, a rede vista pela porta é linear e produz $U''=-R_{Th}I$.}
\end{questao}

\begin{questao}[4]{}
Mostre que os modelos abaixo são equivalentes para qualquer carga e obtenha as relações entre seus parâmetros.
\circuitoq{\begin{circuitikz}[american,scale=.7,transform shape]
\draw (0,0) to[\fonteV,l^={$V_{Th}$}] (0,3.4) to[R,l={$R_{Th}$}] (3.5,3.4)--(4.5,3.4) node[ocirc]{};
\draw (0,0)--(4.5,0) node[ocirc]{}; \node[above] at(2.2,4){Thévenin};
\draw (7,0) to[isource,l^={$I_N$}] (7,3.4)--(10.3,3.4)--(11.3,3.4) node[ocirc]{};
\draw (10.3,3.4) to[R,l_={$R_N$}] (10.3,0); \draw (7,0)--(11.3,0) node[ocirc]{}; \node[above] at(9.1,4){Norton};
\end{circuitikz}}
\resposta{$R_N=R_{Th}$ e $I_N=V_{Th}/R_{Th}$}
\resolucao{Thévenin impõe $U=V_{Th}-R_{Th}I$. Norton impõe $U=I_NR_N-R_NI$. As duas retas coincidem quando intercepto e inclinação são iguais.}
\end{questao}

\begin{questao}[4]{}
Deduza um método de mínimos quadrados para estimar $V_{Th}$ e $R_{Th}$ a partir de $n\ge3$ medições $(I_k,U_k)$.
\resposta{Ajustar a reta $U=V_{Th}-R_{Th}I$; intercepto $=V_{Th}$ e inclinação $=-R_{Th}$}
\resolucao{Use regressão linear: $-R_{Th}=\sum(I_k-\bar I)(U_k-\bar U)/\sum(I_k-\bar I)^2$ e $V_{Th}=\bar U+R_{Th}\bar I$. Resíduos sistemáticos indicam não linearidade.}
\end{questao}

\begin{questao}[4]{}
Projete um ensaio seguro para estimar o Norton de uma fonte de $\SI{60}{\volt}$ cuja corrente de curto prevista é $\SI{120}{\ampere}$, usando instrumentos limitados a $\SI{25}{\ampere}$.
\circuitoq{\begin{circuitikz}[american,scale=.82,transform shape]
\draw[dashed,cinza,rounded corners] (0,0) rectangle (3.3,3.8); \node[cinza] at(1.65,2){fonte};
\draw (3.3,3.1)--(4.8,3.1) coordinate(a); \draw (3.3,.7)--(4.8,.7) coordinate(b); \draw (a) to[R,l={$R_L$},i>^={$I$}] (b);
\node[right,align=left,font=\footnotesize] at(5.7,2){usar dois ou mais\\pontos abaixo de 25 A};
\end{circuitikz}}
\resposta{Exemplo: $R_L=\SI{5.5}{\ohm}$ e $\SI{2.5}{\ohm}$, produzindo aproximadamente $10$ A e $20$ A para $R_{Th}\approx0{,}5\,\Omega$}
\resolucao{Escolhem-se cargas que mantenham a corrente abaixo do limite e suficientemente separadas para evitar mau condicionamento. Ajusta-se a reta $U=V_{Th}-R_{Th}I$ e extrapola-se $I_N=V_{Th}/R_{Th}$ sem curto real.}
\end{questao}

\begin{questao}[4]{}
A rede de duas portas não pode ser representada simultaneamente por dois equivalentes independentes se houver acoplamento. Explique a necessidade de uma matriz de porta.
\circuitoq{\begin{circuitikz}[american,scale=.82,transform shape]
\draw[dashed,cinza,rounded corners] (0,0) rectangle (4.2,4); \node[cinza] at(2.1,2){rede acoplada};
\draw (0,3.2)--(-1.4,3.2) node[ocirc]{} node[left]{A+}; \draw (0,.8)--(-1.4,.8) node[ocirc]{} node[left]{A-};
\draw (4.2,3.2)--(5.6,3.2) node[ocirc]{} node[right]{B+}; \draw (4.2,.8)--(5.6,.8) node[ocirc]{} node[right]{B-};
\end{circuitikz}}
\resposta{É necessário um modelo de dois acessos, por exemplo $\mathbf U=\mathbf V_0-\mathbf R\mathbf I$}
\resolucao{Uma corrente aplicada em A pode alterar a tensão em B e vice-versa. Os termos fora da diagonal da matriz representam esse acoplamento, algo que dois Thévenins independentes não capturam.}
\end{questao}

\begin{questao}[4]{}
Para o equivalente abaixo, maximize a potência na carga sob a restrição $U_L\ge\SI{20}{\volt}$.
\circuitoq{\begin{circuitikz}[american,scale=.88,transform shape]
\draw (0,0) to[\fonteV,l^={$\SI{24}{\volt}$}] (0,3.8) to[R,l={$\SI{8}{\ohm}$}] (4.2,3.8) to[R,l={$R_L$}] (4.2,0)--(0,0);
\draw[<->,Vcol] (5,.4)--(5,3.4) node[midway,right] {$U_L\ge20$ V};
\end{circuitikz}}
\resposta{$R_L=\SI{40}{\ohm}$ no limite; $P_L=\SI{10}{\watt}$}
\resolucao{A restrição exige $R_L\ge40\,\Omega$. Para $R_L>R_{Th}$, a potência decresce com $R_L$, portanto o máximo admissível ocorre no menor valor permitido, 40 $\Omega$.}
\end{questao}

\begin{questao}[4]{}
Aplique uma fonte de teste e determine quando a rede ativa abaixo apresenta $R_{Th}<0$.
\circuitoq{\begin{circuitikz}[american,scale=.78,transform shape]
\coordinate (g) at (3.6,0); \coordinate (a) at (3.6,4);
\draw (a) to[R,l_={$R$}] (g); \draw (0,0) to[cisource,l^={$g_mV_A$}] (0,4)--(a); \draw (0,0)--(g);
\draw (7.2,0) to[isource,l_={$I_t$}] (7.2,4)--(a); \draw (7.2,0)--(g); \node[ground] at(g){};
\end{circuitikz}}
\resposta{$R_{Th}=1/(1/R-g_m)$; é negativo quando $g_m>1/R$}
\resolucao{$I_t=V_A/R-g_mV_A$. Logo $R_{Th}=V_A/I_t=1/(1/R-g_m)$. Se a transcondutância superar a condutância passiva, a inclinação da porta muda de sinal.}
\end{questao}

\begin{questao}[4]{}
Dois divisores entregam a mesma tensão em vazio, mas têm resistências de Thévenin diferentes. Formule um critério de projeto para escolher a porta mais robusta à variação de carga.
\resposta{Preferir a menor $R_{Th}$ para manter a tensão mais próxima de $V_{Th}$ sob carga}
\resolucao{$U_L=V_{Th}R_L/(R_{Th}+R_L)$. Para um mesmo $V_{Th}$ e $R_L$, a queda relativa cresce monotonicamente com $R_{Th}$. Portanto a porta de menor resistência de saída é mais rígida.}
\end{questao}

\begin{questao}[4]{}
Uma caracterização experimental produz pontos que não ficam sobre uma reta. Dê três causas físicas e indique como separar erro de medição de não linearidade real.
\resposta{Aquecimento, limitação de corrente/saturação e elementos não lineares; repetir em níveis menores e com mais pontos}
\resolucao{Erros aleatórios geram resíduos sem padrão; não linearidade tende a produzir curvatura sistemática. Repetir o ensaio com amplitudes menores ajuda a revelar aquecimento, saturação ou mudança de estado.}
\end{questao}

\begin{questao}[4]{}
Uma rede será usada com vinte cargas diferentes. Compare a estratégia de resolver o circuito completo a cada carga com a de obter uma vez $V_{Th}$ e $R_{Th}$.
\resposta{O equivalente exige uma caracterização inicial e depois apenas uma expressão escalar por carga; a análise completa repete o sistema vinte vezes}
\resolucao{O ganho cresce com o número de cargas. Thévenin é especialmente vantajoso quando a rede interna permanece fixa e apenas a carga varia; perde vantagem se forem necessárias grandezas internas da rede.}
\end{questao}
